% Generated draft: requires mathematical and source review.
\documentclass[11pt]{article}
\usepackage[margin=25mm]{geometry}
\usepackage{fontspec}
\setmainfont{Latin Modern Roman}
\usepackage{amsmath,amssymb,mathtools,unicode-math}
\setmathfont{Latin Modern Math}
\usepackage{enumitem,needspace,xurl,hyperref}
\setlength{\parindent}{0pt}
\setlength{\parskip}{5pt}
\setlength{\emergencystretch}{3em}
\newcommand{\N}{\mathbb{N}}
\newcommand{\Z}{\mathbb{Z}}
\newcommand{\Q}{\mathbb{Q}}
\newcommand{\R}{\mathbb{R}}
\newcommand{\C}{\mathbb{C}}
\newcommand{\F}{\mathbb{F}}
\newcommand{\A}{\mathbb{A}}
\newcommand{\PP}{\mathbb{P}}
\newcommand{\E}{\mathbb{E}}
\newcommand{\eps}{\varepsilon}
\newcommand{\dd}{\,\mathrm{d}}
\begin{document}
\section*{UnsolvedMath ID: 30005590}
\textbf{Source problem code:} catalogue entry
\subsection{Definitions and mathematical statement}
The supplied record concerns surface maps satisfying a Hamiltonian-minimal equation obtained from Legendrian area variation. The equation itself, its weak formulation, the domain, the ambient contact manifold and metric, the admissible variations, and any boundary or integrability hypotheses are not supplied. Thus the class quantified over by ``every weak solution'' cannot yet be defined precisely.

For terminology only, if a contact distribution is locally written as \(\ker\alpha\), the Legendrian tangency condition for a smooth surface map \(u\), in the appropriate ambient dimension, is \(u^*\alpha=0\). A branched Legendrian immersion is intended to allow branch points while being a Legendrian immersion elsewhere; the required local model at a branch point is not specified in the record. Neither the precise Schoen--Wolfson cone models nor the Maslov-class construction for their links is supplied.

\textbf{Intended conjecture, with unresolved hypotheses.} Every weak solution of the Hamiltonian-minimal equation arising from Legendrian area variation is a smooth branched Legendrian immersion away from isolated Schoen--Wolfson conical singularities, each having nonzero Maslov class.

Here the exceptional conical singularities are asserted to be isolated; the record does not specify whether accumulation at a boundary or at an end is permitted. It also does not specify the sense in which a singularity is conical. The original excerpt says ``any solution to this equation,'' whereas the reconstructed statement specifies weak solutions. The missing antecedent equation and solution class must be recovered before this can be treated as a fully specified mathematical conjecture.
\subsection{Short English statement}
The proposed regularity statement says that every weak Hamiltonian-minimal Legendrian surface is smooth, allowing branch points, except at isolated Schoen--Wolfson cone singularities with nonzero Maslov class. The supplied record omits the equation and hypotheses needed to specify exactly which surfaces and weak solutions are covered.
\subsection{Sources}
\begin{itemize}
\item Supplied UnsolvedMath local registry record, permanent ID \(30005590\), title: \emph{Regularity of Weak Hamiltonian-Minimal Legendrian Surfaces}; exact source problem code: \texttt{catalogue entry}. Collection: Oberwolfach Reports Open Problems. The proposer is not identified; the supplied proposed year is 2023.
\item Supplied source citation: \emph{Partial Differential Equations (2024)}, Oberwolfach Reports, pp.~1789--1842, DOI \texttt{10.4171/owr/2023/32}. \url{https://doi.org/10.4171/owr/2023/32}. The supplied record describes the reports as documenting workshop discussions and open problems.
\item Supplied literature reference labelled ``Riviere'': \url{https://doi.org/10.1007/s42543-024-00090-y}. No full bibliographic citation or explicit proposer attribution is supplied.
\item Supplied literature reference labelled ``variational theory'': \url{https://arxiv.org/abs/2405.06955}.
\item Supplied literature reference labelled ``original OWR report'': \url{https://doi.org/10.4171/owr/2023/32}.
\item Literature assessment dated exactly \texttt{2026-08-22}: the supplied curation labels the problem open and states that the conjecture remains explicit in the primary paper, while recent variational and partial-regularity results establish important special cases without proving regularity for every weak Hamiltonian-minimal solution. This is supplied dated curation derived from the OWR material, not an independently verified literature assessment; its status and cited claims require independent verification.
\item The registry labels the statement \texttt{corrected\_verified} and says it was reconstructed and checked against the official report. These are supplied provenance assertions, not verification performed for this draft. The original excerpt also appends a passage naming Bozhidar Velichkov, jointly with Guido De Philippis and Luca Spolaor, under the title ``On the fine structure of the two-phase free boundaries''; its relationship to this conjecture is unresolved and it does not establish attribution here.
\end{itemize}
\end{document}
